\documentclass[11pt,a4paper]{report}

\usepackage[utf8]{inputenc}
\usepackage[T1]{fontenc}
\usepackage[brazil]{babel}
\usepackage{lmodern}
\usepackage{amsmath,amssymb}
\usepackage{booktabs}
\usepackage{longtable}
\usepackage{array}
\usepackage{listings}
\usepackage{xcolor}
\usepackage{geometry}
\usepackage{tikz}
\usepackage{fancyhdr}
\usepackage[hidelinks]{hyperref}

\usetikzlibrary{arrows.meta,positioning,shapes.geometric,fit,backgrounds}

\geometry{left=3cm,right=2.7cm,top=2.8cm,bottom=2.8cm}

\definecolor{codebg}{RGB}{247,247,245}
\definecolor{codekw}{RGB}{20,70,140}
\definecolor{codecom}{RGB}{110,110,110}
\definecolor{codestr}{RGB}{140,60,20}
\definecolor{rule}{RGB}{190,190,185}
\definecolor{callout}{RGB}{242,244,248}

\lstset{
  basicstyle=\ttfamily\footnotesize,
  backgroundcolor=\color{codebg},
  keywordstyle=\color{codekw}\bfseries,
  commentstyle=\color{codecom}\itshape,
  stringstyle=\color{codestr},
  showstringspaces=false,
  breaklines=true,
  frame=leftline,
  framerule=1.2pt,
  rulecolor=\color{rule},
  xleftmargin=10pt,
  framexleftmargin=8pt,
  framexrightmargin=4pt,
  aboveskip=10pt,
  belowskip=10pt,
  literate=%
    {á}{{\'a}}1 {é}{{\'e}}1 {í}{{\'i}}1 {ó}{{\'o}}1 {ú}{{\'u}}1
    {â}{{\^a}}1 {ê}{{\^e}}1 {ô}{{\^o}}1 {ã}{{\~a}}1 {õ}{{\~o}}1
    {ç}{{\c c}}1 {à}{{\`a}}1
}

\newcommand{\cpn}{\textsc{CPNCheck}}
\newcommand{\pyra}{\textsc{pyRuleAnalyzer}}
\newcommand{\cpntools}{CPN~Tools}
\newcommand{\code}[1]{\texttt{\small #1}}

\newenvironment{nota}
  {\par\vspace{6pt}\noindent\begin{tikzpicture}\node[fill=callout,rounded corners=2pt,
   inner sep=9pt,text width=\dimexpr\linewidth-20pt\relax,align=justify]\bgroup\small}
  {\egroup;\end{tikzpicture}\par\vspace{6pt}}

\pagestyle{fancy}
\fancyhf{}
\fancyhead[L]{\small\nouppercase{\leftmark}}
\fancyhead[R]{\small\thepage}
\renewcommand{\headrulewidth}{0.4pt}

\title{\vspace{-1.5cm}\Huge\normalfont\scshape CPNCheck\upshape\\[6pt]
\Large Model checking de Redes de Petri Coloridas\\[4pt]
\large fundamentos, implementação e por que confiar no resultado}
\author{Documentação técnica}
\date{\today}

\begin{document}

\maketitle
\thispagestyle{empty}

\vspace{1cm}
\begin{center}
\begin{minipage}{0.82\textwidth}
\small
\noindent\textbf{Sobre este documento.} Ele começa do zero: o que é model
checking, para que serve, e por que a resposta que ele dá é diferente da
resposta que um teste dá. Depois desce até a implementação do \cpn{} módulo
por módulo, e termina com a parte que mais importa para quem vai defender um
resultado --- o que exatamente está garantido, o que não está, e quais
evidências sustentam cada afirmação.

\vspace{6pt}
\noindent Não é preciso conhecer lógica temporal nem Redes de Petri para
começar. Os capítulos~\ref{cap:mc} e~\ref{cap:cpn} constroem os dois conceitos
do início.
\end{minipage}
\end{center}

\vfill
\begin{center}
\small
\url{https://github.com/GrupoCybersegurancaVirtus/cpncheck}
\end{center}

\clearpage
\tableofcontents
\clearpage

% ============================================================================
\chapter{O problema}
% ============================================================================

\section{Um teste responde uma pergunta menor do que parece}

Suponha que você exportou um modelo de \emph{machine learning} --- uma árvore
de decisão, uma floresta aleatória, um \emph{gradient boosting} --- para uma
Rede de Petri Colorida, e quer ter certeza de que a rede se comporta bem. A
primeira coisa que ocorre é testar: rodar a rede com mil entradas e conferir
mil saídas.

Isso é valioso, mas responde a uma pergunta estreita: \emph{``para estas mil
entradas, a rede deu a resposta certa?''}. Não responde:

\begin{itemize}
  \item a rede sempre \emph{para}, ou existe entrada que a faz girar para
        sempre?
  \item existe algum estado em que ela trava sem produzir resposta?
  \item existem duas folhas da mesma árvore que podem estar habilitadas ao
        mesmo tempo --- ou seja, a rede é realmente determinística?
  \item existe alguma execução em que o resultado é depositado duas vezes?
\end{itemize}

Nenhuma dessas perguntas é sobre uma entrada específica. Todas são sobre
\emph{todos os comportamentos possíveis} do modelo. Teste é amostragem:
executa alguns caminhos e conclui por indução informal. Model checking é
exaustão: enumera \emph{todos} os estados alcançáveis e decide a pergunta
sobre o conjunto inteiro.

\begin{nota}
\textbf{A distinção que vale memorizar.} Um teste que passa diz ``não
encontrei problema nestes casos''. Um model checker que responde ``verdadeiro''
diz ``não existe problema, ponto'' --- dentro do modelo e do escopo que ele
explorou, que é exatamente o que este documento vai tornar preciso.
\end{nota}

\section{O que é preciso para verificar}

Model checking precisa de três coisas, e a qualidade da resposta depende das
três:

\begin{enumerate}
  \item \textbf{Um modelo formal} do sistema --- aqui, a Rede de Petri
        Colorida lida do arquivo \code{.cpn};
  \item \textbf{Uma propriedade formal}, escrita numa lógica com significado
        matemático preciso --- aqui, CTL;
  \item \textbf{Um algoritmo} que decide se o modelo satisfaz a propriedade,
        e que, quando não satisfaz, exibe \emph{por quê}.
\end{enumerate}

O terceiro ponto é o que separa model checking de uma prova de teorema comum:
quando a resposta é ``falso'', a ferramenta devolve um \textbf{contraexemplo}
--- uma execução concreta que viola a propriedade. É a diferença entre um
médico dizer ``algo está errado'' e apontar exatamente onde.

% ============================================================================
\chapter{Model checking do começo}
\label{cap:mc}
% ============================================================================

\section{O modelo: sistemas de transição}

Tudo em model checking parte de um objeto simples. Um \textbf{sistema de
transição} (ou \emph{estrutura de Kripke}) é uma quádrupla
\[
  K = (S,\ s_0,\ R,\ L)
\]
onde:

\begin{itemize}
  \item $S$ é um conjunto de \textbf{estados};
  \item $s_0 \in S$ é o \textbf{estado inicial};
  \item $R \subseteq S \times S$ é a \textbf{relação de transição}: $(s,s')
        \in R$ significa ``de $s$ o sistema pode ir para $s'$'';
  \item $L$ rotula cada estado com as \textbf{proposições atômicas} que valem
        nele --- fatos elementares como ``o place \emph{Prediction} tem token''.
\end{itemize}

Um \textbf{caminho} é uma sequência $s_0 \to s_1 \to s_2 \to \cdots$ seguindo
$R$. Um caminho é \textbf{maximal} quando não pode ser estendido: ou é
infinito, ou termina num estado sem sucessores.

\begin{center}
\begin{tikzpicture}[
  node distance=16mm and 20mm,
  st/.style={circle,draw,minimum size=8mm,inner sep=1pt,font=\small},
  every edge/.style={draw,-{Stealth[length=2mm]}}]
  \node[st,label=below:{\footnotesize inicial}] (a) {$s_0$};
  \node[st,right=of a] (b) {$s_1$};
  \node[st,right=of b] (c) {$s_2$};
  \node[st,above right=8mm and 20mm of b] (d) {$s_3$};
  \node[st,right=of c,fill=black!8] (e) {$s_4$};
  \draw (a) edge (b);
  \draw (b) edge (c);
  \draw (b) edge (d);
  \draw (c) edge (e);
  \draw (d) edge[bend left=10] (e);
  \node[below=1mm of e,font=\footnotesize] {sem sucessores};
\end{tikzpicture}
\end{center}

Nesse exemplo há dois caminhos maximais: $s_0 s_1 s_2 s_4$ e $s_0 s_1 s_3 s_4$.
Ambos terminam --- $s_4$ é o que chamaremos de \emph{marcação morta}.

\section{A propriedade: lógica temporal CTL}

\textbf{CTL} (\emph{Computation Tree Logic}) descreve como proposições se
comportam \emph{ao longo do tempo} e \emph{ao longo dos caminhos possíveis}.
Cada operador CTL é a combinação de duas escolhas:

\begin{center}
\begin{tabular}{@{}ll@{}}
\toprule
\textbf{Quantificador de caminho} & \textbf{Operador temporal} \\
\midrule
\textbf{A} --- em \emph{todos} os caminhos & \textbf{X} --- no próximo estado \\
\textbf{E} --- em \emph{algum} caminho     & \textbf{F} --- em \emph{algum} estado futuro \\
                                            & \textbf{G} --- em \emph{todos} os estados futuros \\
                                            & \textbf{U} --- até que \\
\bottomrule
\end{tabular}
\end{center}

Combinando, obtêm-se os operadores usados na prática. Esta é a tabela mais
importante do documento:

\begin{center}
\begin{tabular}{@{}llp{7.6cm}@{}}
\toprule
\textbf{Fórmula} & \textbf{Nome} & \textbf{Leitura} \\
\midrule
$AG\,\varphi$ & invariante, segurança
  & $\varphi$ vale em \emph{toda} marcação alcançável. ``Nada de ruim acontece.'' \\[4pt]
$AF\,\varphi$ & inevitabilidade, vivacidade
  & \emph{toda} execução maximal chega a um estado onde $\varphi$ vale. ``Algo de bom acontece, sempre.'' \\[4pt]
$EF\,\varphi$ & alcançabilidade
  & \emph{existe} execução que chega a $\varphi$. ``É possível chegar lá.'' \\[4pt]
$E[\varphi\ U\ \psi]$ & até que (existencial)
  & existe execução em que $\varphi$ vale continuamente até que $\psi$ passe a valer. \\
\bottomrule
\end{tabular}
\end{center}

\subsection*{Por que essas quatro bastam}

As demais se derivam. Em particular:
\[
  \neg EF\,\varphi \;\equiv\; AG\,\neg\varphi,
  \qquad
  \neg E[\varphi\ U\ \psi] \;=\; \text{``$\psi$ nunca acontece antes de $\varphi$''}
\]

A segunda é a forma das propriedades de \textbf{precedência}: ``o estágio $m$
nunca é concluído antes do estágio $m{-}1$''. Guarde essa leitura; ela
reaparece no catálogo (propriedades C1 e C2).

\subsection*{Segurança e vivacidade são diferentes na prática}

A distinção não é acadêmica, e ela determina como é o contraexemplo:

\begin{itemize}
  \item Um \textbf{invariante} ($AG$) é violado por \emph{um único estado ruim}.
        O contraexemplo é finito: a sequência mais curta do início até esse
        estado.
  \item Uma \textbf{vivacidade} ($AF$) é violada por \emph{uma execução
        inteira} que nunca chega ao objetivo. Se essa execução for infinita,
        \emph{não existe contraexemplo finito} --- a menos que se exiba a
        parte que se repete. Essa é a forma \emph{laço} (\emph{lasso}), e o
        capítulo~\ref{cap:cex} mostra como o \cpn{} a produz.
\end{itemize}

\section{O algoritmo: pontos fixos}

CTL é decidida sem enumerar caminhos --- o que seria inviável, já que há
infinitos. Em vez disso, calcula-se para cada fórmula $\varphi$ o
\textbf{conjunto de satisfação} $\mathrm{Sat}(\varphi) \subseteq S$: os estados
onde ela vale. A resposta final é simplesmente $s_0 \in \mathrm{Sat}(\varphi)$.

Os conjuntos são calculados por \textbf{pontos fixos}, de dentro para fora:

\paragraph{$EF$ e $E[\,U\,]$: alcançabilidade para trás.}
$\mathrm{Sat}(EF\,\varphi)$ é o menor conjunto que contém
$\mathrm{Sat}(\varphi)$ e é fechado por predecessores. Algoritmo: parta dos
estados onde $\varphi$ vale e caminhe \emph{para trás} pelas arestas até
saturar.

\paragraph{$AF$: ``todo sucessor já está''.}
$\mathrm{Sat}(AF\,\varphi)$ é o menor conjunto $Z$ tal que
$Z = \mathrm{Sat}(\varphi) \cup \{s : \text{todo sucessor de } s \text{ está em } Z\}$.
Implementa-se com um contador por estado: quando o último sucessor de $s$ entra
no conjunto, $s$ entra também.

\paragraph{$AG$: o complemento de uma alcançabilidade.}
$AG\,\varphi$ vale em $s$ exatamente quando \emph{nenhum} estado alcançável a
partir de $s$ viola $\varphi$. Então: calcule os estados ruins, calcule quem
alcança um estado ruim (para trás), e tome o complemento.

\begin{nota}
\textbf{Custo.} Cada operador custa uma varredura do grafo, então o total é
$O(|\varphi| \times (|S| + |R|))$ --- \emph{linear} no tamanho do grafo. O
gargalo nunca é a lógica; é o tamanho de $S$.
\end{nota}

\section{O ponto delicado: marcações mortas}
\label{sec:dead}

CTL de livro-texto assume que a relação $R$ é \textbf{total}: todo estado tem
ao menos um sucessor, e portanto todo caminho é infinito. Mas os modelos aqui
são feitos para \emph{parar}: a rede classifica a entrada e termina numa
\textbf{marcação morta} (nenhuma transição habilitada).

Isso obriga a decidir o que $AF\,\varphi$ significa num caminho que acaba.
Existem duas convenções, e elas dão respostas \emph{diferentes}:

\begin{center}
\begin{tabular}{@{}p{3.3cm}p{10.2cm}@{}}
\toprule
\textbf{Convenção} & \textbf{O que diz numa marcação morta sem $\varphi$} \\
\midrule
\textbf{livro-texto}
  & A marcação morta ganha um laço sobre si mesma (Baier \& Katoen, seção 3.1).
    O caminho vira infinito, nunca alcança $\varphi$, e portanto
    \textbf{falsifica} $AF\,\varphi$. \\[5pt]
\textbf{ASK-CTL} (\cpntools)
  & O \code{EV} do ASK-CTL é \emph{vacuamente verdadeiro} numa marcação sem
    sucessores. Assim \code{EV($\varphi$)} significa na verdade
    $AF(\varphi \vee \mathit{morta})$: uma rede que trava antes de decidir
    \textbf{satisfaz} \code{EV(PRED)}. \\
\bottomrule
\end{tabular}
\end{center}

A diferença aparece exatamente quando existe uma marcação morta alcançável sem
$\varphi$ --- ou seja, no caso em que você mais quer ser avisado.

\begin{nota}
\textbf{Como o \cpn{} resolve.} Ele adota a leitura de livro-texto (mais
severa) e, ao gerar o script para o \cpntools{}, codifica $AF\,\varphi$ como
\[
  \texttt{EV}(\varphi)\ \wedge\ \neg\,\texttt{E}[\neg\varphi\ \texttt{U}\ (\mathit{morta} \wedge \neg\varphi)]
\]
para que o \cpntools{} responda à \emph{mesma} pergunta. Além disso, reporta
as duas leituras (chaves \code{A3\_ev}, \code{B2\_ev}, \code{C3a\_ev},
\code{D2a\_ev}), de modo que a diferença fique visível em vez de escondida.
\end{nota}

\section{O limite real: explosão de estados}

O número de estados cresce exponencialmente com o paralelismo do modelo. Para
os modelos de \emph{gradient boosting} gerados pelo \pyra{}, com $K$ classes e
$M$ estágios por classe, o grafo de ocorrências tem exatamente
\[
  (3M+3)^K + 2 \ \text{nós} \qquad\text{e}\qquad K(3M+2)(3M+3)^{K-1} + 2 \ \text{arcos},
\]
porque os canais de classe intercalam livremente. Um modelo de três classes com
cem estágios está fora do alcance de qualquer ferramenta explícita --- inclusive
o \cpntools{}.

Isso não é defeito da implementação; é a natureza do problema. O que uma
ferramenta honesta faz é \textbf{recusar-se a responder} quando não terminou a
exploração, e é o que o \cpn{} faz: acima do orçamento (\code{max\_nodes},
200.000 por padrão) as propriedades de espaço de estados ficam \emph{sem
veredicto}, enquanto a análise estrutural e o teste de conformidade continuam
valendo.

% ============================================================================
\chapter{Redes de Petri Coloridas}
\label{cap:cpn}
% ============================================================================

\section{Rede de Petri}

Uma Rede de Petri tem dois tipos de nó:

\begin{itemize}
  \item \textbf{Lugares} (\emph{places}), desenhados como elipses, que guardam
        \textbf{fichas} (\emph{tokens});
  \item \textbf{Transições} (\emph{transitions}), desenhadas como retângulos,
        que consomem fichas de alguns lugares e produzem em outros.
\end{itemize}

Os \textbf{arcos} ligam lugares a transições e vice-versa. O estado do sistema
é a \textbf{marcação}: quantas fichas há em cada lugar. Uma transição está
\textbf{habilitada} quando há fichas suficientes nos lugares de entrada;
quando \textbf{dispara}, consome-as e produz as de saída.

\begin{center}
\begin{tikzpicture}[
  pl/.style={ellipse,draw,minimum width=17mm,minimum height=8mm,font=\small},
  tr/.style={rectangle,draw,fill=black!6,minimum width=6mm,minimum height=12mm,font=\small},
  every edge/.style={draw,-{Stealth[length=2mm]}},
  node distance=14mm]
  \node[pl,label=above:{\footnotesize $\bullet$}] (p1) {Entrada};
  \node[tr,right=of p1] (t) {$t$};
  \node[pl,right=of t] (p2) {Saída};
  \draw (p1) edge (t);
  \draw (t) edge (p2);
  \node[below=6mm of t,font=\footnotesize,align=center]
    {guarda: \code{[\#f0 x <= 3.0]}};
\end{tikzpicture}
\end{center}

\section{O que ``colorida'' acrescenta}

Numa Rede de Petri comum as fichas são indistinguíveis. Numa
\textbf{Rede de Petri Colorida} (CPN) cada ficha carrega um \textbf{valor} ---
uma ``cor''. Isso traz quatro elementos:

\begin{description}
  \item[Conjuntos de cores (\emph{colour sets}).] O tipo dos valores:
        \code{INT}, \code{REAL}, \code{BOOL}, produtos, registros, listas. No
        modelo de uma floresta, por exemplo, \code{SAMPLE} é um registro de
        seis reais.
  \item[Variáveis.] Declaradas com um conjunto de cores:
        \code{var x: SAMPLE;}
  \item[Inscrições de arco.] Expressões em CPN ML. No arco de entrada é um
        \emph{padrão} que casa com a ficha; no de saída é uma expressão que
        constrói a ficha produzida.
  \item[Guardas.] Condições booleanas na transição, como
        \code{[\#f0 x <= 3.0 andalso \#f1 x > 1.5]}.
\end{description}

\subsection*{Elementos de ligação}

Aqui está o conceito que mais confunde e que o \cpn{} trata explicitamente.
Uma transição por si só não é uma ação: uma ação é um \textbf{elemento de
ligação} (\emph{binding element}) --- a transição \emph{mais} uma atribuição
de valores às suas variáveis.

As variáveis ganham valor de dois modos:

\begin{enumerate}
  \item \textbf{Por casamento de padrão}: a variável aparece num arco de
        entrada, e seu valor vem da ficha consumida.
  \item \textbf{Por enumeração sobre o conjunto de cores}: a variável aparece
        só na guarda ou num arco de saída. Nesse caso o \cpntools{} percorre
        \emph{todos} os valores do conjunto de cores dela.
\end{enumerate}

O segundo caso é fácil de esquecer e caro de ignorar: uma ferramenta que não
o implementa está analisando uma rede \emph{menor} do que a desenhada --- e
não avisa. O \cpn{} implementa a enumeração e, quando o conjunto de cores é
infinito (um \code{REAL}, por exemplo), \textbf{recusa a rede nomeando a
variável} em vez de seguir em frente.

\section{Hierarquia}

Redes grandes se organizam em \textbf{páginas}. Uma \emph{transição de
substituição} não dispara: ela representa uma subpágina inteira. Os lugares da
página de cima (\emph{sockets}) são fundidos com os lugares de borda da
subpágina (\emph{ports}).

Nos modelos gerados pelo \pyra{}: a página de topo tem \code{Input} e
\code{Prediction}; cada árvore ou canal de classe é uma subpágina.

Para verificar, o \cpn{} \textbf{achata} a hierarquia: funde cada par
port/socket num único lugar e produz uma rede plana equivalente. Nada se perde
--- a hierarquia é organização visual, não semântica.

\section{O grafo de ocorrências}

Aqui os dois capítulos se encontram. O \textbf{grafo de ocorrências}
(\emph{occurrence graph}, ou espaço de estados) de uma CPN é construído assim:

\begin{itemize}
  \item cada \textbf{nó} é uma marcação alcançável;
  \item cada \textbf{arco} é a ocorrência de um elemento de ligação.
\end{itemize}

\textbf{Esse grafo é exatamente a estrutura de Kripke do capítulo~\ref{cap:mc}.}
As marcações são os estados $S$, as ocorrências são a relação $R$, a marcação
inicial é $s_0$, e as proposições atômicas são fatos sobre as fichas (``o lugar
\emph{Prediction} tem ao menos uma ficha''). Por isso CTL se aplica diretamente.

Uma sutileza importante: as marcações carregam os \textbf{valores reais} das
fichas. Não há abstração. Duas marcações que diferem no valor de uma única
ficha são \emph{dois nós distintos}. Isso torna a verificação exata para a
entrada dada, e é também por isso que a verificação é feita \emph{por entrada}
--- ponto ao qual voltaremos na seção~\ref{sec:escopo}.

\section{Três conceitos derivados do grafo}

\begin{description}
  \item[Marcação morta.] Nó sem sucessores: a rede parou. Numa rede correta,
        toda marcação morta contém uma predição.
  \item[Componente fortemente conexa (SCC).] Conjunto maximal de nós
        mutuamente alcançáveis. A decomposição em SCCs (algoritmo de Tarjan)
        revela ciclos.
  \item[\emph{Home marking}.] Marcação alcançável a partir de \emph{qualquer}
        marcação alcançável. Existe exatamente quando o grafo de SCCs tem uma
        \emph{única} componente terminal (Jensen \& Kristensen). É o que o
        \code{ListHomeMarkings} do \cpntools{} calcula.
\end{description}

% ============================================================================
\chapter{Para que serve isso num modelo de \emph{machine learning}}
% ============================================================================

\section{O que se está verificando --- e o que não se está}
\label{sec:escopo}

Vale ser muito claro, porque é aqui que se cometem exageros.

\begin{nota}
O \cpn{} verifica a \textbf{rede exportada}, lida de volta do arquivo
\code{.cpn}. Ele \textbf{não} verifica que o modelo de ML é acurado, justo,
robusto a ataque adversarial ou adequado ao problema. Ele verifica que o
\emph{artefato formal} que você vai publicar, embarcar ou citar se comporta
como um classificador deve se comportar --- e que ele computa a mesma coisa
que o classificador de origem nas entradas testadas.
\end{nota}

Dito isso, as garantias obtidas são reais e não-triviais:

\begin{itemize}
  \item \textbf{Determinismo.} Duas folhas da mesma árvore nunca estão
        habilitadas simultaneamente --- a árvore não é ambígua.
  \item \textbf{Terminação.} A rede sempre para; não existe ciclo alcançável.
  \item \textbf{Ausência de travamento.} Toda parada é uma parada \emph{com
        resposta}.
  \item \textbf{Saída única e válida.} Nunca duas predições; nunca um rótulo
        fora do domínio declarado.
  \item \textbf{Disjunção das guardas.} As folhas de uma árvore particionam o
        espaço de entrada --- e esta vale para \emph{toda} entrada, não só as
        exploradas.
  \item \textbf{Conformidade.} Nas entradas testadas, a rede prediz o mesmo
        que o classificador.
\end{itemize}

\section{A escada de escopo}

As quatro técnicas do \cpn{} têm alcances \emph{diferentes}, e confundir isso
é o erro mais fácil de cometer ao relatar resultados:

\begin{center}
\begin{tabular}{@{}llp{7.2cm}@{}}
\toprule
\textbf{Técnica} & \textbf{Escopo} & \textbf{Um veredicto ``verdadeiro'' significa} \\
\midrule
\code{ctl} & uma entrada
  & Para \emph{esta} marcação inicial, a propriedade vale em todos os
    comportamentos possíveis. Outra entrada é outro grafo. \\[4pt]
\code{scc} & uma entrada
  & Idem; é uma pergunta sobre a estrutura do mesmo grafo. \\[4pt]
\code{structural} & \textbf{toda} entrada
  & Vale para o espaço de entrada inteiro, sem explorar nada. É o veredicto
    mais forte que a ferramenta produz. \\[4pt]
\code{testing} & as entradas dadas
  & Equivalência funcional nas entradas fornecidas. É teste, não prova. \\
\bottomrule
\end{tabular}
\end{center}

\begin{nota}
\textbf{Como relatar sem exagerar.} ``As propriedades A1--A8 e B1--B3 foram
verificadas para $n$ entradas do conjunto de teste; B4 (disjunção de guardas)
vale para todo o espaço de entrada; PC foi testada em $N$ amostras.'' Isso é
exato. Dizer ``o modelo foi formalmente verificado'' sem qualificação não é.
\end{nota}

% ============================================================================
\chapter{\cpn{} por dentro}
% ============================================================================

\section{Arquitetura}

O \cpn{} tem cerca de 9.000 linhas de Python e \textbf{nenhuma dependência de
terceiros} --- só a biblioteca padrão. As camadas vão do arquivo até o
veredicto:

\begin{center}
\begin{tikzpicture}[
  box/.style={rectangle,draw,rounded corners=2pt,text width=74mm,
              inner xsep=3mm,inner ysep=1.7mm,font=\small,fill=black!3,
              align=left},
  side/.style={rectangle,draw,rounded corners=2pt,text width=38mm,
               inner xsep=2mm,inner ysep=1.7mm,font=\small,fill=black!8,
               align=center},
  lbl/.style={font=\scriptsize\itshape,text=black!55,text width=38mm,
              align=center},
  node distance=2.6mm]
  \node[box] (net) {\code{net.py} \hfill lê o XML do \cpntools{}};
  \node[box,below=of net] (ml) {\code{cpnml.py} \hfill lexer, parser, compilador};
  \node[box,below=of ml] (cs) {\code{colourset.py} \hfill cores e domínios};
  \node[box,below=of cs] (col) {\code{coloured.py} \hfill achata, liga, dispara};
  \node[box,below=of col] (ss) {\code{statespace.py} \hfill grafo, SCCs};
  \node[box,below=of ss] (ctl) {\code{ctl.py} \hfill pontos fixos CTL};
  \node[box,below=of ctl] (f) {\code{formula.py} \hfill fórmulas como dados};
  \node[box,below=of f] (ch) {\code{checker.py} \hfill orquestra e reporta};

  \node[side,right=14mm of col] (ask) {\code{askctl.py}};
  \node[lbl,above=1.5mm of ask] {a mesma fórmula,\\para o \cpntools{}};

  \node[side,right=14mm of f] (prof) {\code{profile.py}\\[1pt]\code{profiles/}};
  \node[lbl,below=1.5mm of prof] {o que o chamador sabe\\sobre \emph{suas} redes};

  \draw[-{Stealth[length=2mm]}] (prof.west) -- (f.east);
  \draw[-{Stealth[length=2mm]},rounded corners=3pt]
        (f.east) -- ++(6mm,0) -- ++(0,26mm) -- (ask.south);
\end{tikzpicture}
\end{center}

\section{Camada por camada}

\subsection{\code{net.py} --- ler o arquivo}

Faz \emph{só} o \emph{parsing} do XML que o \cpntools{} grava: páginas,
lugares com suas marcações iniciais, transições com guardas, arcos com
inscrições, transições de substituição com seus pares port/socket.

São 96 linhas, e não sabem nada sobre o que a rede significa. Isso é
deliberado: a mesma leitura serve para qualquer rede.

\subsection{\code{cpnml.py} --- entender as expressões}

Guardas e inscrições são escritas em \textbf{CPN ML}, um subconjunto de
Standard ML. Uma versão anterior as traduzia com expressões regulares --- o
que funciona para o que um gerador emite e falha para o que uma pessoa
escreve.

Hoje é um \emph{parser} de verdade: analisador léxico, analisador sintático por
escalada de precedência, e um compilador que transforma a árvore sintática em
\emph{closures} Python. Cobre:

\begin{center}
\begin{tabular}{@{}lp{9.6cm}@{}}
\toprule
expressões & \code{if/then/else}, \code{let/in/end}, \code{case/of}, \code{fn}, aplicação de função \\
operadores & a tabela do SML com precedência e associatividade: \code{orelse}, \code{andalso}, comparações, \code{::}, \code{@}, \code{+ - \^{}}, \code{* / div mod}, \code{\~{}} e \code{not} unários \\
valores & inteiros, reais (inclusive \code{\~{}} e \code{E\~{}}), strings, booleanos, unit, tuplas, listas, registros, \emph{multisets} \\
padrões & variáveis, curinga, literais, tuplas, listas, \code{::}, registros, anotações de tipo \\
declarações & \code{val} e \code{fun}, incluindo funções multi-cláusula e recursão \\
\bottomrule
\end{tabular}
\end{center}

\begin{nota}
\textbf{Consequência que vale destacar: a rede se auto-descreve.} As funções
\code{vadd}, \code{vsum}, \code{argmax}, \code{allzero} e \code{decide} usadas
pelos modelos de floresta \emph{não} são funções Python embutidas no
verificador. Elas são lidas dos blocos \code{<ml>} do próprio arquivo
\code{.cpn}, compiladas e executadas. Nada sobre votação suave ou
\emph{boosting} está dentro da ferramenta.
\end{nota}

A compilação para \emph{closures} (em vez de interpretação nó a nó) mantém o
disparo rápido: a árvore é percorrida uma única vez, na compilação.

\subsection{\code{colourset.py} --- os domínios}

Lê os conjuntos de cores e as variáveis dos elementos XML estruturados
(\code{<color>}, \code{<var>}) e decide quais são \emph{finitos}:

\begin{center}
\begin{tabular}{@{}lp{10.2cm}@{}}
\toprule
\textbf{Finitos} & \code{unit}, \code{bool}, enumerações, faixas de índice, inteiros com limites, e produtos/registros/subconjuntos/\emph{aliases} construídos a partir desses \\
\textbf{Infinitos} & \code{int} sem limites, \code{real}, \code{string}, listas \\
\bottomrule
\end{tabular}
\end{center}

Um domínio infinito é reportado como ``desconhecido'', nunca chutado.

\subsection{\code{coloured.py} --- a rede que dispara}

Achata a hierarquia (fundindo ports e sockets), compila cada guarda e
inscrição, e implementa as duas operações centrais:

\begin{description}
  \item[\code{enabled(marcação)}] enumera os elementos de ligação habilitados:
        casa os padrões dos arcos de entrada contra as fichas disponíveis e,
        em seguida, percorre os domínios das variáveis livres.
  \item[\code{fire(marcação, transição, ligação)}] calcula a marcação sucessora:
        remove as fichas consumidas, avalia as inscrições de saída, deposita
        as produzidas.
\end{description}

As fichas são valores Python escolhidos para serem \emph{hasheáveis e
ordenáveis}: registro é tupla na ordem dos campos declarados, lista é tupla,
\emph{unit} é a tupla vazia. Assim duas marcações são o mesmo nó exatamente
quando são iguais --- o que é o que permite usar um dicionário para
reconhecer marcações já visitadas.

\subsection{\code{statespace.py} --- o grafo}

Busca em largura a partir da marcação inicial, com orçamento de nós. Para cada
nó guarda o \emph{pai} na árvore de busca --- incluindo a \textbf{ligação}
usada, sem a qual um traço não poderia ser redisparado. Em seguida, decomposição
em SCCs por Tarjan iterativo, de onde saem as componentes terminais e os
\emph{home markings}.

\subsection{\code{ctl.py} --- os pontos fixos}

Três funções, exatamente as do capítulo~\ref{cap:mc}:

\begin{center}
\begin{tabular}{@{}lp{9.4cm}@{}}
\toprule
\code{\_eu} & $E[\varphi\ U\ \psi]$ por alcançabilidade para trás \\
\code{\_af} & $AF\,\varphi$ por contagem de sucessores, com o parâmetro
              \code{semantics} escolhendo entre a leitura de livro-texto e a do
              ASK-CTL (seção~\ref{sec:dead}) \\
\code{\_af\_lasso} & extrai o \emph{laço} quando $AF$ falha numa execução
              infinita (capítulo~\ref{cap:cex}) \\
\bottomrule
\end{tabular}
\end{center}

\subsection{\code{formula.py} --- propriedades como dados}

Este é o módulo que dá flexibilidade à ferramenta. Uma propriedade não é código
Python escrito contra uma rede específica; é uma \emph{árvore de objetos}:

\begin{center}
\begin{tabular}{@{}lp{9.2cm}@{}}
\toprule
\textbf{Proposições atômicas} & \code{Dead}, \code{Marked}, \code{AllMarked}, \code{MaxTokens}, \code{TokensIn}, \code{Enabled}, \code{EnabledCount} \\
\textbf{Conectivos} & \code{Not}, \code{And}, \code{Or}, \code{Implies} \\
\textbf{Temporais} & \code{AG}, \code{AF}, \code{EF}, \code{EU} \\
\textbf{Quantificador} & \code{ForEach} --- ``para toda árvore\ldots'', ``para todo canal\ldots'' \\
\bottomrule
\end{tabular}
\end{center}

Cada fórmula faz duas coisas além de ser verdadeira ou falsa:

\begin{enumerate}
  \item \textbf{explica-se} quando falha --- o \code{AG} devolve a sequência
        mais curta até a marcação violadora e pergunta ao operando \emph{por
        que aquela marcação é ruim};
  \item \textbf{renderiza-se em ASK-CTL} --- a \emph{mesma} árvore vira a
        consulta para o \cpntools{}.
\end{enumerate}

\begin{nota}
\textbf{Por que isso importa para a confiança.} Antes, a verificação em Python
e o script ASK-CTL eram escritos à mão, lado a lado. Nada além de cuidado
impedia que divergissem --- e uma divergência silenciosa significaria comparar
respostas para perguntas diferentes. Hoje há uma única árvore e dois
\emph{backends}, e a divergência é impossível por construção.
\end{nota}

\subsection{\code{profile.py} e \code{profiles/} --- onde entra o domínio}

O verificador não sabe que existe um lugar chamado \code{Prediction}. Quem
sabe é o \textbf{perfil}, que responde:

\begin{center}
\begin{tabular}{@{}lp{8.2cm}@{}}
\toprule
\code{matches} / \code{family} & esta rede é uma das minhas, e como eu a chamo \\
\code{resolve\_groups} & os domínios que o \code{ForEach} percorre \\
\code{ctl\_formulas} & as propriedades temporais, como fórmulas \\
\code{scc\_properties} & as respondidas pela decomposição em SCCs \\
\code{static\_checks} & vereditos que dispensam grafo de ocorrências \\
\code{conformance} & comparação com uma implementação de referência \\
\code{label\_domain} & os valores de saída admissíveis \\
\code{as\_input\_token} / \code{input\_places} & como outra entrada entra na rede \\
\code{mutation\_operators} & as falhas que as propriedades devem pegar \\
\code{unavailable} & propriedades que a rede tem mas esta execução não decide \\
\code{catalog} & os descritores sob os quais os vereditos são reportados \\
\bottomrule
\end{tabular}
\end{center}

O perfil \code{ml\_trees} é o exemplo completo: ele reconhece as redes do
\pyra{} e declara as vinte propriedades do catálogo.

\begin{nota}
\textbf{A distinção entre ``não se aplica'' e ``não pôde ser decidida''} foi
um defeito real encontrado tarde. Uma propriedade que a rede simplesmente não
tem (agregação numa rede que não agrega) é omitida e nunca mencionada --- certo.
Mas uma propriedade que a rede \emph{tem} e que não pôde ser decidida (o
domínio de rótulos de uma árvore de decisão, que não está escrito em lugar
algum legível) estava sendo omitida também --- o que se lê como se tivesse
passado. Hoje ela é reportada como \code{SKIP} com o motivo.
\end{nota}

\section{As quatro técnicas}

\begin{center}
\begin{tabular}{@{}lp{10.4cm}@{}}
\toprule
\code{ctl}
  & Model checking explícito sobre o grafo de ocorrências. Valores reais de
    ficha, sem abstração. \\[4pt]
\code{scc}
  & \emph{Home markings} pela decomposição em componentes fortemente conexas.
    Não é uma fórmula CTL. \\[4pt]
\code{structural}
  & Análise das guardas que são conjunções de limites. Uma guarda
    \code{\#f0 x <= 3.0 andalso \#f1 x > 1.5} descreve uma \emph{caixa}; duas
    guardas podem valer juntas exatamente quando as caixas se intersectam.
    Vale para \emph{toda} entrada. \\[4pt]
\code{testing}
  & Conformidade com uma implementação de referência nas entradas dadas. Segue
    uma sequência maximal por entrada, o que é barato o bastante para um
    conjunto de teste inteiro. \\
\bottomrule
\end{tabular}
\end{center}

\begin{nota}
A técnica \code{testing} só é significativa \emph{ao lado} do model checking,
nunca no lugar dele: ela segue \emph{uma} execução maximal por entrada, e que
essa execução dê ``o'' resultado da rede depende de confluência --- que é
justamente o que as propriedades de determinismo estabelecem.
\end{nota}

\section{Exatidão da análise estrutural}

Um detalhe que mudou durante o desenvolvimento e que vale registrar, porque é
um caso de resposta errada dada em silêncio.

A versão original tratava os extremos fechados como se fossem abertos,
mapeando $\ge$ em $>$ e $<$ em $\le$, e testava a sobreposição com
$\max(\text{inferiores}) < \min(\text{superiores})$.

Considere então as guardas $x \ge 3$ e $x \le 3$, que compartilham o ponto
$x = 3$. O teste dava $3 < 3$, ou seja \emph{disjuntas} --- um
\textbf{falso negativo}, que é exatamente o erro que importa: dizer ``não há
sobreposição'' quando há.

A versão atual rastreia se cada extremo é aberto ou fechado e decide
exatamente:

\begin{center}
\begin{tabular}{@{}lll@{}}
\toprule
$y \ge 3$ \ \emph{vs} \ $y \le 3$ & sobrepõem em $\{3\}$ & \checkmark \\
$y > 3$   \ \emph{vs} \ $y \le 3$ & disjuntas            & \checkmark \\
$y \ge 3$ \ \emph{vs} \ $y < 3$   & disjuntas            & \checkmark \\
\bottomrule
\end{tabular}
\end{center}

Nenhum veredicto sobre as redes geradas mudou --- o exportador só emite
\code{<=} e \code{>}, onde as duas versões concordam. Mas para uma rede escrita
à mão era uma resposta errada esperando para acontecer.

\section{O catálogo de propriedades}

As vinte propriedades das redes do \pyra{}. Note a coluna \emph{Técnica}: ela
determina o escopo do veredicto, conforme a seção~\ref{sec:escopo}.

\begin{center}
\footnotesize
\begin{longtable}{@{}llp{5.6cm}l@{}}
\toprule
\textbf{ID} & \textbf{Propriedade} & \textbf{Fórmula} & \textbf{Técnica} \\
\midrule
\endhead
\multicolumn{4}{@{}l}{\textit{Comuns às três famílias}}\\[2pt]
A1  & Terminação                  & $AF\ \mathit{morta}$                              & ctl \\
A2  & Sem travamento espúrio      & $\neg EF(\mathit{morta} \wedge \neg \mathit{pred})$ & ctl \\
A3  & Decisão inevitável          & $AF\ \mathit{pred}$                               & ctl \\
A4  & Predição recuperável        & $AG\ EF\ \mathit{pred}$                           & ctl \\
A5  & Saída única                 & $AG\ |\mathit{Prediction}| \le 1$                 & ctl \\
A6  & Rótulo válido               & $AG\ \mathit{Prediction} \subseteq L$             & ctl \\
A7  & Segurança (\emph{safeness}) & $AG\ \forall p:\ |p| \le 1$                       & ctl \\
A8  & \emph{Home marking}         & uma única componente terminal                     & scc \\[4pt]
\multicolumn{4}{@{}l}{\textit{Por árvore}}\\[2pt]
B1  & Determinismo de folha       & $\forall T:\ AG\ |EN(\mathit{folhas}_T)| \le 1$   & ctl \\
B2  & Seleção de folha inevitável & $\forall T:\ AF\ EN(\mathit{folhas}_T)$           & ctl \\
B3  & Saída única por árvore      & $\forall T:\ AG\ |\mathit{out}_T| \le 1$          & ctl \\
B4  & Disjunção das guardas       & $\forall T, i \ne j:\ \mathit{caixa}_i \cap \mathit{caixa}_j = \emptyset$ & \textbf{structural} \\[4pt]
\multicolumn{4}{@{}l}{\textit{Gradient boosting}}\\[2pt]
C1  & Precedência de estágio      & $\forall m:\ \neg E[\neg \mathit{acc}_{m-1}\ U\ \mathit{acc}_m]$ & ctl \\
C2  & Sem escore prematuro        & $\neg E[\neg \mathit{acc}_M\ U\ \mathit{score}]$ & ctl \\
C3a & Disparo de decisão inevitável & $AF\ EN(\mathit{decide})$                       & ctl \\
C3b & Determinismo da decisão     & $AG\ |EN(\mathit{decide})| \le 1$                 & ctl \\[4pt]
\multicolumn{4}{@{}l}{\textit{Floresta aleatória}}\\[2pt]
D1  & Voto completo               & $AG(EN(\mathit{vote}) \rightarrow \mathit{todos\ votaram})$ & ctl \\
D2a & Agregação inevitável        & $AF\ EN(\mathit{vote})$                           & ctl \\
D2b & Agregação única             & $AG\ |\mathit{Prediction}| \le 1$                 & ctl \\[4pt]
\multicolumn{4}{@{}l}{\textit{Conformidade}}\\[2pt]
PC  & Consistência da predição    & $\mathit{rede}(x) = \mathit{modelo}(x)$           & \textbf{testing} \\
\bottomrule
\end{longtable}
\end{center}

\noindent onde $EN(G)$ vale numa marcação com arco de saída cuja transição está
em $G$, e $|EN(G)|$ conta esses arcos. As propriedades C1 e C2 são
quantificadas também sobre os canais de classe $k$, omitido acima por
brevidade: $\mathit{acc}_m$ é o acumulador do estágio $m$ do canal, e $M$ é o
último estágio.

% ============================================================================
\chapter{Contraexemplos}
\label{cap:cex}
% ============================================================================

\section{Por que um veredicto sozinho não basta}

``A propriedade A3 é falsa'' é uma informação pobre. Ela não diz se o problema
é do modelo, da exportação ou da propriedade; não permite reproduzir; não
permite corrigir. Um model checker sério devolve a \textbf{evidência}.

No \cpn{} o objeto \code{Counterexample} carrega a sequência de ocorrências,
a \textbf{ligação de cada passo}, a marcação final e a \textbf{entrada} que a
violação exige.

\section{As cinco formas}

\begin{center}
\begin{tabular}{@{}lp{10.2cm}@{}}
\toprule
\code{finite}
  & Uma propriedade de segurança falhou: existe marcação alcançável que a
    viola, e o prefixo é a sequência mais curta até ela. \\[4pt]
\code{lasso}
  & Uma propriedade de vivacidade falhou numa execução infinita: um prefixo
    até um ciclo que a rede pode repetir para sempre sem nunca chegar ao
    objetivo. \textbf{Sem o ciclo, este caso não tem testemunha finita
    alguma.} \\[4pt]
\code{witness}
  & Uma falha que não é uma sequência --- por exemplo, duas componentes
    terminais sem sucessor comum. \\[4pt]
\code{structural}
  & Um defeito da rede como escrita, independente de execução: duas guardas
    que podem valer juntas. Carrega um \emph{ponto} do espaço de entrada onde
    isso acontece. \\[4pt]
\code{conformance}
  & A rede e a referência discordam numa entrada concreta. \\
\bottomrule
\end{tabular}
\end{center}

\section{O laço, e por que ele precisou de um algoritmo}

Quando $AF\,\varphi$ falha, há dois motivos possíveis:

\begin{enumerate}
  \item a execução \emph{para} antes de $\varphi$ --- e então o contraexemplo é
        um caminho finito até uma marcação morta fora de $\varphi$;
  \item a execução \emph{continua para sempre} sem alcançar $\varphi$.
\end{enumerate}

No segundo caso não existe caminho finito que sirva de prova. A versão anterior
da ferramenta apenas informava ``existe um ciclo alcançável que evita
\ldots'', sem nada para examinar.

O algoritmo \code{\_af\_lasso} extrai o ciclo, e é correto por construção. O
argumento:

\begin{nota}
Seja $S$ o conjunto dos nós que \emph{não} satisfazem $AF\,\varphi$. Nenhum
nó de $S$ satisfaz $\varphi$. Além disso, todo nó de $S$ que tem sucessores
tem \emph{um sucessor em $S$} --- pois se todos os sucessores satisfizessem
$AF\,\varphi$, o próprio nó satisfaria. Logo, uma caminhada a partir da
marcação inicial dentro de $S$ só pode parar numa marcação morta. Quando
nenhuma marcação morta de $S$ é alcançável --- que é exatamente o caso em que
a busca por caminho finito não achou nada --- a caminhada \textbf{tem} de
repetir um nó, e a repetição é o ciclo.
\end{nota}

Na prática o traço sai assim:

\begin{lstlisting}[language={}]
A3  [lasso] a predicao nunca acontece
    sequence: Distribute<x = (-1.36, -1.91, ...)>
              -> Init<x = (...)>
              -> Feed_1<sc = -0.043, x = (...)>
              then forever [GBDT1T1_Rule2_Class1<x = (...)>]
    cycle_length: 1
\end{lstlisting}

\section{O certificado: \code{replay()}}

Esta é a parte que sustenta a confiança na ferramenta.

Um contraexemplo é uma sequência finita de elementos de ligação. O método
\code{replay()} \textbf{redispara} essa sequência a partir da marcação inicial
usando a \emph{regra de disparo da rede} --- não a busca que produziu o traço
--- e confere que ela termina onde o contraexemplo afirma. Para um laço,
confere também que o ciclo retorna ao ponto de entrada.

\begin{nota}
\textbf{Por que isso é forte.} A busca (pontos fixos, BFS, Tarjan) e a regra de
disparo são código independente. Um erro nos pontos fixos \emph{não consegue
forjar} um traço que redispara. Um contraexemplo que passa no \code{replay()}
é um contraexemplo que a ferramenta não inventou.

A suíte de testes exige que \emph{toda} sequência de \emph{todo} mutante
redispare --- e exige que uma sequência adulterada (invertida ou truncada)
\textbf{não} redispare. Sem esse segundo teste, o certificado não certificaria
nada.
\end{nota}

\section{De achado estrutural para entrada concreta}

A disjunção de guardas (B4) vale ou falha para \emph{toda} entrada, então uma
falha não nomeia execução nenhuma. Mas a interseção de duas guardas em caixa é
uma caixa, e um ponto dela é uma \textbf{entrada} onde a sobreposição de fato
acontece.

O \cpn{} calcula esse ponto e o devolve como \code{cex.input}. Realimentando-o:

\begin{lstlisting}[language=Python]
resultado = check_cpn("modelo.cpn", class_labels=[0, 1])
# B4 FAIL  |  B1 PASS   (a entrada original nao alcanca o par)

x = resultado.properties["B4"].cex.input     # (1.31, 0, 0, 1.117, 0, 0)

de_novo = check_cpn("modelo.cpn", samples=[x], class_labels=[0, 1])
# B4 FAIL  |  B1 FAIL   -- e no MESMO par de folhas
\end{lstlisting}

Isso fecha o ciclo entre a prova estrutural (toda entrada) e a prova de estado
(uma entrada): o achado deixa de ser uma afirmação e passa a ser
\emph{reprodutível}. É a ideia de refinamento guiado por contraexemplo
(CEGIS) aplicada aqui.

% ============================================================================
\chapter{Por que confiar no resultado}
% ============================================================================

\section{O problema de fundo}

Um verificador é um programa, e programas têm defeitos. Como saber que um
``VERIFIED'' não é apenas um \emph{bug} que responde ``verdadeiro'' para tudo?

Duas evidências independentes sustentam a resposta, e é importante que sejam
\emph{duas}, porque cada uma sozinha é insuficiente.

\section{Primeira evidência: sensibilidade (mutantes)}

De uma rede correta, o \cpn{} deriva \textbf{dezesseis} redes defeituosas ---
cada uma uma pequena edição de guarda, inscrição ou arco, escolhida para
quebrar \emph{exatamente} uma propriedade nomeada:

\begin{center}
\footnotesize
\begin{tabular}{@{}llp{7.0cm}@{}}
\toprule
\textbf{Mutante} & \textbf{Deve falhar} & \textbf{Edição} \\
\midrule
\code{cycle}          & A1      & folha viva devolve a ficha de entrada \\
\code{deadlock}       & A2, A3  & produtor da predição perde o arco de saída \\
\code{dead\_branch}   & A4, A8  & folha dormente sempre habilitada, sem produzir \\
\code{double\_output} & A5, A7  & predição depositada duas vezes \\
\code{bad\_label}     & A6      & predição substituída pelo rótulo 7 \\
\code{overlap}        & B1, B4  & folha dormente habilitada junto com a viva \\
\code{hidden\_overlap}& B4      & duas folhas dormentes com a mesma guarda \\
\code{no\_leaf}       & B2      & guarda da folha viva vira falsa \\
\code{tree\_double}   & B3      & árvore emite o resultado duas vezes \\
\code{skip\_stage}    & C1      & \code{Feed\_2} lê \code{Acc0} \\
\code{early\_score}   & C2      & \code{Finalize} lê \code{Acc}$_{M-1}$ \\
\code{no\_decision}   & C3a     & guarda do \code{Decide} vivo vira falsa \\
\code{two\_decisions} & C3b     & guarda de um \code{Decide} dormente vira verdadeira \\
\code{partial\_vote}  & D1      & \code{Vote} deixa de consumir o voto da árvore 0 \\
\code{no\_vote}       & D2a     & guarda do \code{Vote} vira falsa \\
\code{wrong\_value}   & PC      & rede prediz outra classe válida \\
\bottomrule
\end{tabular}
\end{center}

\begin{nota}
\textbf{O par \code{overlap} / \code{hidden\_overlap} merece atenção.} O
primeiro habilita uma folha dormente ao lado da viva: a entrada do arquivo
alcança as duas, e tanto B1 (por entrada) quanto B4 (estrutural) falham. O
segundo dá a mesma guarda a duas folhas que a entrada \emph{nunca alcança}:
B1 \textbf{passa}, e só B4 falha. Esse par demonstra que as duas técnicas têm
escopos genuinamente diferentes --- não é retórica.

Do mesmo modo, \code{wrong\_value} produz uma rede que satisfaz \emph{todas}
as propriedades CTL e falha apenas na conformidade: comportar-se bem e computar
a coisa certa são perguntas distintas.
\end{nota}

\section{Segunda evidência: concordância com o \cpntools{}}

O \cpn{} executa o simulador do \cpntools{} 4.0.1 de forma \emph{headless} ---
através do Access/CPN, entrando na ferramenta de espaço de estados exatamente
como o \emph{Enter State Space} da interface faz, calculando os grafos de
ocorrências e de SCCs e avaliando consultas ASK-CTL --- e compara cada
estatística e cada veredicto.

\begin{center}
\begin{tabular}{@{}lr@{}}
\toprule
Redes comparadas (duas corretas + todos os mutantes de cada) & \textbf{28} \\
Redes em que \emph{todos} os valores coincidem & \textbf{28} \\
Vereditos \emph{falsos} confirmados pelo \cpntools{} & \textbf{81} \\
Discordâncias & \textbf{0} \\
\bottomrule
\end{tabular}
\end{center}

\begin{nota}
\textbf{Por que os 81 são a parte que importa.} Concordância apenas nas redes
corretas também seria alcançada por dois verificadores que sempre respondem
``verdadeiro''. São os vereditos \emph{falsos} --- confirmados independentemente
pelo \cpntools{} --- que mostram que ambos os motores detectam as mesmas
violações nos mesmos lugares.
\end{nota}

Essa comparação foi refeita \emph{depois} de o motor de avaliação do CPN ML ser
inteiramente substituído (expressões regulares $\to$ \emph{parser} real):
mesmas 28 redes, mesmos 81 vereditos falsos, nenhuma discordância. Nenhum
veredicto se moveu.

\section{Terceira evidência: os contraexemplos redisparam}

Como descrito no capítulo~\ref{cap:cex}, toda sequência produzida para todo
mutante é redisparada com a regra de disparo da rede e tem de reproduzir a
violação. É uma verificação da busca por código que não é a busca.

\section{O que \emph{não} está garantido}

Esta seção é tão importante quanto as anteriores.

\begin{center}
\begin{tabular}{@{}p{5.4cm}p{8.2cm}@{}}
\toprule
\textbf{Limitação} & \textbf{Consequência prática} \\
\midrule
Verificação \code{ctl} é \textbf{por entrada}
  & Verificar 3 amostras não diz nada sobre a quarta. Só B4 (estrutural) vale
    para todo o espaço. \\[4pt]
Explosão de estados
  & Acima de \code{max\_nodes} as propriedades de espaço de estados ficam
    \emph{sem veredicto} --- e ``sem veredicto'' nunca deve ser lido como
    ``aprovado''. Quatro reduções opcionais adiam esse limite (conjuntos
    teimosos, \emph{sweep-line}, equivalência e simetria); ver o documento
    sobre explosão de estados. \\[4pt]
PC é teste
  & Conformidade nas entradas dadas, não prova de equivalência funcional. \\[4pt]
O fragmento CPN ML tem limites
  & Tempo real (\code{@++}), valores aleatórios, canais síncronos e código com
    efeitos colaterais não são suportados (redes temporizadas com tempo
    inteiro, prioridades e conjuntos de fusão são). O que está fora levanta
    \code{UnsupportedNet} em vez de ser mal interpretado. \\[4pt]
Variável livre sobre domínio infinito
  & A rede é \textbf{recusada}, nomeando a variável. Prosseguir verificaria uma
    rede menor do que a desenhada. \\[4pt]
Verifica a rede, não o modelo
  & Acurácia, justiça e robustez do classificador são outras perguntas. \\
\bottomrule
\end{tabular}
\end{center}

\begin{nota}
\textbf{O princípio de projeto por trás de tudo isso}: quando a ferramenta não
sabe, ela \emph{diz que não sabe}. Nunca chuta, nunca omite em silêncio. Três
defeitos encontrados no desenvolvimento eram exatamente violações desse
princípio --- uma propriedade sumindo do relatório em vez de aparecer como
\code{SKIP}; uma entrada sendo ignorada sem aviso; uma sobreposição de guardas
reportada como disjunta. Os três foram corrigidos, e cada um virou teste.
\end{nota}

% ============================================================================
\chapter{Uso prático}
% ============================================================================

\section{Linha de comando}

\begin{lstlisting}[language=bash]
cpncheck check modelo.cpn --labels 0,1,2     # verifica e imprime o relatorio
cpncheck check modelo.cpn --trace            # + contraexemplos completos
cpncheck check modelo.cpn --json out.json    # vereditos em JSON
cpncheck check modelo.cpn --latex            # tabela LaTeX
cpncheck askctl modelo.cpn -o askctl.sml     # script para o CPN Tools
cpncheck mutants modelo.cpn mutantes/        # um mutante por propriedade
cpncheck oracle modelo.cpn                   # roda o CPN Tools e compara
\end{lstlisting}

\section{Python}

\begin{lstlisting}[language=Python]
from cpncheck import check_cpn

resultado = check_cpn("modelo.cpn", class_labels=[0, 1, 2])
print(resultado.report())
assert resultado.passed

for r in resultado.failures:
    print(r.id, r.detail)          # o contraexemplo, renderizado
    cex = r.cex
    cex.kind                       # finite | lasso | witness | structural | conformance
    cex.sequence(bindings=True)    # com as ligacoes de cada passo
    cex.input                      # a entrada que a violacao exige
    cex.replay(checker.cnet)       # True: ele realmente faz o que diz
\end{lstlisting}

\section{Como ler o relatório}

\begin{lstlisting}[language={},basicstyle=\ttfamily\scriptsize]
Model checking: gbdt_binary.cpn  [Gradient Boosting Decision Trees]
=======================================================================
  occurrence graph : 14 nodes, 13 arcs
  SCC graph        : 14 nodes, 13 arcs, 1 dead marking(s), 1 home marking(s)
  inputs           : 1 model checked, 0 tested   time: 0.00s
  semantics        : CTL over the occurrence graph, dead markings with
                     a self-loop

  CTL model checking on the occurrence graph
  ---------------------------------------------------------------------
  A1   Termination                  AF dead                         PASS
  ...
  16 passed, 0 failed, 0 skipped -- VERIFIED
\end{lstlisting}

Três coisas a observar:

\begin{enumerate}
  \item A linha \code{inputs} diz \emph{quantas entradas} foram verificadas.
        Um \code{VERIFIED} com \code{1 model checked} é uma afirmação sobre
        uma entrada.
  \item A linha \code{semantics} declara a convenção de marcação morta em uso.
  \item \code{skipped} não é \code{passed}. Uma propriedade pulada traz o
        motivo ao lado.
\end{enumerate}

\section{Verificando redes que não são de \emph{machine learning}}

As propriedades são dados, não código. Um perfil próprio permite verificar
qualquer rede:

\begin{lstlisting}[language=Python]
import collections
from cpncheck import CPNModelChecker
from cpncheck.formula import AF, AG, Dead, MaxTokens
from cpncheck.profile import Profile
from cpncheck.selectors import Places

class MeuPerfil(Profile):
    name = "minhas redes"
    catalog = meus_descritores

    def matches(self, net):
        return any(p["name"] == "Done" for p in net.places)

    def ctl_formulas(self, net, groups, labels):
        return collections.OrderedDict([
            ("T1", AF(Dead().phrased("o fim"))),
            ("T2", AG(MaxTokens(Places(name="Done"), 1))),
        ])

print(CPNModelChecker("minha.cpn", profile=MeuPerfil()).check().report())
\end{lstlisting}

% ============================================================================
\chapter*{Glossário}
\addcontentsline{toc}{chapter}{Glossário}
% ============================================================================

\begin{description}\setlength{\itemsep}{3pt}
  \item[Marcação] O estado da rede: quantas fichas, e com que valores, há em
        cada lugar.
  \item[Elemento de ligação] Uma transição \emph{mais} uma atribuição de
        valores às suas variáveis. É a unidade de ação.
  \item[Grafo de ocorrências] O grafo de todas as marcações alcançáveis; a
        estrutura de Kripke da rede.
  \item[Marcação morta] Marcação sem transições habilitadas: a rede parou.
  \item[\emph{Home marking}] Marcação alcançável a partir de qualquer marcação
        alcançável.
  \item[SCC] Componente fortemente conexa: conjunto maximal de nós mutuamente
        alcançáveis.
  \item[CTL] Lógica temporal de árvores de computação; $AG$, $AF$, $EF$,
        $E[\cdot\,U\,\cdot]$.
  \item[Segurança / vivacidade] ``Nada de ruim acontece'' ($AG$) \emph{vs.}
        ``algo de bom acontece'' ($AF$).
  \item[Contraexemplo] Evidência concreta de que uma propriedade não vale.
  \item[Laço (\emph{lasso})] Contraexemplo com prefixo + ciclo, a forma
        necessária quando a violação é uma execução infinita.
  \item[Mutante] Rede derivada de uma correta por uma edição que quebra uma
        propriedade nomeada; serve para provar que o verificador consegue
        responder ``falso''.
  \item[Perfil] O objeto que informa ao verificador o que as redes de um
        domínio significam.
  \item[ASK-CTL] A biblioteca de model checking do \cpntools{}.
  \item[CPN ML] O subconjunto de Standard ML usado em guardas e inscrições.
\end{description}

\vfill
\begin{center}\small
\rule{0.4\textwidth}{0.4pt}\\[6pt]
\cpn{} --- MIT\\
\url{https://github.com/GrupoCybersegurancaVirtus/cpncheck}
\end{center}

\end{document}
